\chapter{Linear Dynamics}\label{linear:ch:dynamics}

\begin{center}\itshape
The motion of a machine with its governor consists in general of a uniform motion, combined with a disturbance which may be
expressed as the sum of several component motions. These components may be of four different kinds.
\upshape(James Clerk Maxwell, \emph{On Governors}, 1868)
\end{center}
\noindent\emph{Use Appendix~\ref{linear:ch:eigen}, earlier matrix norm and positivity results, and Volume~3's Perron--Frobenius theorem. This optional appendix applies them to stability, iteration and control.}

Watt's centrifugal governor used spinning weighted arms to close a steam valve when an engine ran too fast. Some engines instead hunted: their speed rose and fell in growing oscillations. Tuning by trial and error did not explain why a feedback device could amplify the disturbance it was meant to remove.

Maxwell's 1868 paper \emph{On Governors} linearised the motion near equilibrium. Components of the disturbance had the form $e^{\lambda t}$, so decay required negative real parts. He found the condition for his cubic model and asked for a general criterion. Routh answered in 1877; Hurwitz found the criterion independently in 1895 in connection with turbine design.

The matrix tests are simple to state: powers decay when every eigenvalue has modulus below one, and exponentials decay when every eigenvalue has negative real part. Their proofs also explain why a nonnormal matrix can grow before decaying. Lyapunov gives a quadratic energy certificate; Chebyshev and heavy-ball iteration change the polynomial acting on the error; Kalman's rank condition describes which states an input can reach.

\begin{think}
For $A=\begin{pmatrix}0.9&10\\0&0.9\end{pmatrix}$, both eigenvalues lie inside the unit circle. Compute its first few powers and find the upper-right entry in $A^k$. How can that entry become much larger than one while all powers eventually tend to zero?
\end{think}

\section{Powers of a matrix and Gelfand's formula}

Use positive matrix size and write $\rho(A)=\max|\lambda|$. An induced norm satisfies submultiplicativity and $\rho(A)\le\|A\|$, since it can be tested on an eigenvector. A suitable change of coordinates makes this bound nearly exact.

\begin{lemma}[a norm adapted to $A$]\label{linear:lem:adaptednorm}
For every complex square $A$ and $\varepsilon>0$, some invertible $S$ satisfies $\|S^{-1}AS\|_2\le\rho(A)+\varepsilon$. The vector norm $\|x\|_S=\|S^{-1}x\|_2$ therefore has induced bound $\|A\|_S\le\rho(A)+\varepsilon$.
\end{lemma}
\begin{proof}
Schur gives $A=UTU^*$ with $T=\Lambda+N$ upper triangular and $N$ strictly upper triangular. For $D=\operatorname{diag}(1,\delta,\ldots,\delta^{n-1})$, the entry in position $i,j$ of $D^{-1}TD$ is $t_{ij}\delta^{j-i}$. At $0<\delta\le1$, its strictly upper part has Frobenius norm at most $\delta\|N\|_F$, hence operator norm at most that. Thus $S=UD$ gives bound $\rho(A)+\delta\|N\|_F$; choose $\delta$ small enough. The last assertion follows by substituting $S^{-1}x$ in that operator bound.
\end{proof}

\begin{theorem}[convergence of powers]\label{linear:thm:powers}
Powers $A^k$ tend to zero exactly when $\rho(A)<1$. For every $\varepsilon>0$ and every fixed matrix norm, there is $C_\varepsilon$ with $\|A^k\|\le C_\varepsilon(\rho(A)+\varepsilon)^k$ for all $k\ge0$.
\end{theorem}
\begin{proof}
Conjugating powers gives $\|A^k\|_2\le\|S\|_2\|S^{-1}\|_2(\rho+\varepsilon)^k$. Finite-dimensional norm equivalence proves the bound in any matrix norm. If $\rho<1$, take $\rho+\varepsilon<1$. If an eigenvalue has modulus at least one, its eigenvector satisfies $A^kv=\lambda^kv$, which cannot tend to zero. This proves the converse.
\end{proof}

\begin{theorem}[Gelfand's formula, 1941]\label{linear:thm:gelfand}
For any norm on the matrix space, $\rho(A)=\lim_{k\to\infty}\|A^k\|^{1/k}$.
\end{theorem}
\begin{proof}
For the Euclidean operator norm, an eigenvector gives $\rho(A)^k\le\|A^k\|_2$. The preceding upper bound makes the limit superior of its $k$th root at most $\rho+\varepsilon$. Let $\varepsilon\downarrow0$. Norm-equivalence constants disappear after taking $k$th roots, giving every norm.
\end{proof}

The root rate leaves room for polynomial factors. If $N=E_{12}$, then for $k\ge1$,
\[(\lambda I+cN)^k=\lambda^kI+kc\lambda^{k-1}N,\]
with the zeroth power identity handled separately. At $\lambda=9/10,c=10$, the off-diagonal entry is $10k(9/10)^{k-1}$. Comparing consecutive terms gives ratio $9(k+1)/(10k)$, so the maximum occurs at $k=9,10$ and is about $38.7$. The exact formula gives transient amplification despite eventual decay. The shared root rate in Figure~\ref{linear:fig:transient} is $0.9$; the nonnormal matrix also carries its factor $k$. A normal matrix instead has $\|A^k\|_2=\rho(A)^k$. Such transient amplification can matter before nonlinear terms become important, for example in hydrodynamic stability.

\begin{figure}[ht]
\centering
\fig{linear-ch12-transient}
\caption{$\|A^k\|_2$ for $A=\bigl(\begin{smallmatrix}0.9&10\\0&0.9\end{smallmatrix}\bigr)$ (charcoal, log scale) and for the normal matrix
$0.9I$ with the same eigenvalues (rust). Both decay like $0.9^k$ in the end, but the non-normal one first grows by a factor of almost
$40$.}
\label{linear:fig:transient}
\end{figure}

\section{Differential equations and stability}

The unique solution of $x'=Ax$, $x(0)=x_0$, is $e^{tA}x_0$. Differentiation verifies it, and differentiating $e^{-tA}x(t)$ proves uniqueness. On an eigenvector it grows by $e^{\lambda t}$. Put $\alpha(A)=\max\Re\lambda$; call $A$ \emph{stable} or \emph{Hurwitz} when this is negative.

\begin{theorem}[stability]\label{linear:thm:stability}
For every $\varepsilon>0$, a constant $C_\varepsilon$ satisfies $\|e^{tA}\|\le C_\varepsilon e^{(\alpha(A)+\varepsilon)t}$ for all $t\ge0$. Every solution tends to zero exactly when $A$ is stable, and then decay is exponential.
\end{theorem}
\begin{proof}
Use the triangular scaling from the adapted-norm lemma to write $S^{-1}AS=\Lambda+N_\delta$, where $\|N_\delta\|_2\le\varepsilon$. For $y=S^{-1}x$, differentiate its squared norm:
\[\frac{\dd}{\dd t}\|y\|^2=2\Re\,y^*(\Lambda+N_\delta)y\le2(\alpha+\varepsilon)\|y\|^2.
\]
Multiplication by $e^{-2(\alpha+\varepsilon)t}$ makes it non-increasing. Return through $S$ and use norm equivalence to get the estimate. Negative $\alpha$ permits a negative exponent. Conversely an eigenvector with non-negative real eigenvalue part gives a nondecaying solution, excluding convergence for every initial state.
\end{proof}

The planar examples in Figure~\ref{linear:fig:phase} include nodes, spirals, centres and saddles. Here stability means asymptotic attraction to zero. A centre has bounded closed orbits and is stable in the usual Lyapunov sense of small disturbances staying small, but it is not attracting and is not Hurwitz. Real two-dimensional Hurwitz stability is equivalent to negative trace and positive determinant, as Exercise~\ref{linear:ex:hurwitz2} proves.

\begin{figure}[ht]
\centering
\fig{linear-ch12-phase}
\caption{Phase portraits of $x'=Ax$ in the plane. Stable node and stable spiral: every trajectory tends to $0$. Centre: closed
orbits, neither stable nor unstable in the sense of Theorem~\ref{linear:thm:stability}. Saddle: trajectories come in along the stable
eigenvector and leave along the unstable one (rust).}
\label{linear:fig:phase}
\end{figure}

\section{Lyapunov's method}

Lyapunov's 1892 thesis sought a decreasing energy rather than an explicit solution. For linear systems a quadratic form is enough, and its coefficients satisfy a linear equation.

\portrait{lyapunov}{Aleksandr Lyapunov (1857--1918). Photographer unknown, published by 1908. Public domain, via Wikimedia Commons.}

\begin{theorem}[Lyapunov, 1892]\label{linear:thm:lyapunov}
For complex square $A$, the following are equivalent:
\begin{enumerate}[label=(\roman*)]
\item $A$ is stable;
\item for every $Q\succ0$ there exists $P\succ0$ with $A^*P+PA=-Q$;
\item some $P\succ0$ satisfies $A^*P+PA\prec0$.
\end{enumerate}
When stable, the equation has the unique solution $P=\int_0^\infty e^{tA^*}Qe^{tA}\dd t$.
\end{theorem}
\begin{proof}
(i) implies (ii) because exponential decay makes the integral norm convergent. It is Hermitian and positive definite: for nonzero $x$, its integrand quadratic form is positive at zero and on a nearby interval. Integrating the derivative of $e^{tA^*}Qe^{tA}$ gives $A^*P+PA=-Q$, since the far endpoint vanishes. If $R$ is the difference of two solutions, $e^{tA^*}Re^{tA}$ has zero derivative and equals $R$; exponential decay also makes it tend to zero, so $R=0$.

(ii) implies (iii) by any positive $Q$. For (iii), put $Q=-(A^*P+PA)\succ0$. Along a solution the energy $V=x^*Px$ satisfies
\[V'=-x^*Qx\le-\frac{\lambda_{\min}(Q)}{\lambda_{\max}(P)}V.
\]
Thus $V$ decays exponentially, and $V\ge\lambda_{\min}(P)\|x\|^2$ gives decay of every state. The preceding stability theorem proves (i).
\end{proof}

Trajectories enter smaller energy ellipsoids because the energy derivative is strictly negative away from zero. Solving for $P$ provides a stability certificate without first locating eigenvalues. For a nonlinear system with this stable linearisation, differentiability makes the remaining local terms smaller than the strict quadratic decrease; this is Lyapunov's local linearisation principle. The discrete counterpart is the Stein equation $P-A^*PA=Q$, treated in Exercise~\ref{linear:ex:stein}.

\section{The speed of iterative methods}

Let real symmetric $A\succ0$ have spectrum in $[m,L]$, $0<m\le L$, and solve $Ax=b$, equivalently minimise $f(x)=x^{\mathsf T}Ax/2-b^{\mathsf T}x$. Set $x_\star=A^{-1}b$ and $e_k=x_k-x_\star$. Write $\kappa=L/m$ for the spectral-interval ratio; when the endpoints are the extremal eigenvalues it is the Euclidean condition number.

Gradient descent $x_{k+1}=x_k-\eta(Ax_k-b)$ gives $e_{k+1}=(I-\eta A)e_k$. Symmetry gives the exact operator factor $\rho(I-\eta A)=\max_{\lambda\in\operatorname{spec}A}|1-\eta\lambda|$. The best uniform bound over the interval uses $\eta=2/(m+L)$ and factor $(\kappa-1)/(\kappa+1)$ (Exercise~\ref{linear:ex:gdstep}). For large $\kappa$, a fixed reduction therefore takes order $\kappa$ steps.

More generally suppose $x_k\in x_0+\operatorname{span}(r_0,Ar_0,\ldots,A^{k-1}r_0)$, where $r_0=Ax_0-b=Ae_0$. Then $e_k=p_k(A)e_0$ with $\deg p_k\le k$ and $p_k(0)=1$. This affine Krylov condition is the class for which the normalisation follows; arbitrary linear combinations of iterates need not preserve it. The spectral bound is
\[\|e_k\|_2\le\max_{\lambda\in[m,L]}|p_k(\lambda)|\,\|e_0\|_2.
\]
Chebyshev's cosine polynomials solve the resulting interval minimax problem.

\begin{theorem}[Chebyshev acceleration]\label{linear:thm:chebyshevaccel}
For $0<m<L$, put $\ell(\lambda)=(L+m-2\lambda)/(L-m)$, $\xi_0=\ell(0)=(\kappa+1)/(\kappa-1)$ and $p_k=T_k\circ\ell/T_k(\xi_0)$. Then $p_k(0)=1$ and
\[\max_{[m,L]}|p_k|=\frac1{T_k(\xi_0)}\le2\left(\frac{\sqrt\kappa-1}{\sqrt\kappa+1}\right)^k.
\]
No degree-at-most-$k$ polynomial equal to one at zero has a smaller maximum. For $m=L$, the degree-zero optimum is one and every degree bound $k\ge1$ has optimum zero, attained by $1-\lambda/m$; the displayed divided formula is replaced by this singleton case.
\end{theorem}
\begin{proof}
For $m<L$, the cosine identity gives $|T_k|\le1$ on $[-1,1]$, attained at an endpoint. On $x=\cosh u>1$, the same recurrence gives $T_k(x)=\cosh(ku)\ge e^{ku}/2$. Since $\xi_0+\sqrt{\xi_0^2-1}=(\sqrt\kappa+1)/(\sqrt\kappa-1)$, the bound follows.

For $k\ge1$, $p_k$ takes alternating extreme values at the $k+1$ transformed cosine nodes. If a competitor had a strictly smaller maximum, their difference would have a root between every consecutive pair, giving $k$ roots in $(m,L)$, and also a root at zero. Degree at most $k$ makes this impossible. At $k=0$ the only admissible constant is one. When $m=L$, a linear polynomial can vanish at that one positive point while staying one at zero, proving the separate conclusion.
\end{proof}

The resulting bound needs order $\sqrt\kappa$ steps for a fixed reduction at large $\kappa$. Its three-term recurrence gives an iteration with one matrix-vector multiplication per step. In exact arithmetic conjugate gradients selects the best error polynomial for the actual spectrum and initial error in the $A$ norm, without knowing the interval endpoints; this further algorithmic result is due to Hestenes--Stiefel.

Polyak's 1964 heavy-ball iteration keeps a fixed previous-step correction:
\[x_{k+1}=x_k-\alpha(Ax_k-b)+\beta(x_k-x_{k-1}).\]
\begin{theorem}[heavy ball on quadratics]\label{linear:thm:heavyball}
Take $\alpha=4/(\sqrt L+\sqrt m)^2$ and $\beta=((\sqrt\kappa-1)/(\sqrt\kappa+1))^2$. The paired error iteration has spectral radius $\sqrt\beta$. For every $\varepsilon>0$, $\|e_k\|\le C_\varepsilon(\sqrt\beta+\varepsilon)^k$, with the constant depending on the initial errors. This includes $m=L$, where $\beta=0$ and the paired map is nilpotent.
\end{theorem}
\begin{proof}
In an eigenbasis, each mode pair has matrix
\[M_\lambda=\begin{pmatrix}1+\beta-\alpha\lambda&-\beta\\1&0\end{pmatrix}.
\]
Set $q=(\sqrt L-\sqrt m)/(\sqrt L+\sqrt m)=\sqrt\beta$. Its trace runs affinely from $2q$ at $m$ to $-2q$ at $L$, so its quadratic roots are conjugate or repeated with product $q^2$, both of modulus $q$. For $q=0$ the sole mode matrix is $\begin{pmatrix}0&0\\1&0\end{pmatrix}$ and its square is zero. Taking the block diagonal collection gives the radius in all cases. The power theorem supplies the error bound.
\end{proof}

Figure~\ref{linear:fig:rates} compares the interval bounds at $\kappa=100$. The heavy-ball endpoint modes have repeated roots and nontrivial Jordan blocks when $m<L$, so a polynomial prefactor can accompany the asymptotic rate. Exercise~\ref{linear:ex:heavyjordan} makes it explicit. Nesterov's 1983 acceleration extends the square-root condition-number scale beyond quadratics to smooth strongly convex objectives; that broader result requires a different proof.

\begin{figure}[ht]
\centering
\fig{linear-ch12-rates}
\caption{Worst-case error over the spectrum $[1,100]$ ($\kappa=100$): gradient descent with the optimal step (rate $\frac{99}{101}$), the
heavy ball (rate $\frac9{11}$, with an initial hump), and the Chebyshev bound $2(\frac9{11})^k$.}
\label{linear:fig:rates}
\end{figure}

\section{Controllability}

An input changes the question from where a system goes to where we can send it. Write $x'=Ax+Bu$, with $B\in\C^{n\times p}$ and a continuous input $u:[0,T]\to\C^p$. From $x(0)=0$, variation of constants gives
\[x(T)=\int_0^T e^{(T-s)A}Bu(s)\dd s.
\]
Indeed, differentiate $e^{-tA}x(t)$ and integrate. The system is \emph{controllable} if every target can be reached. Kalman's 1960 criterion reduces this question to rank. His filtering work from the same year later helped guide the Apollo missions.

\begin{theorem}[Kalman's rank condition]\label{linear:thm:kalman}
For each $T>0$, the states reachable from zero with continuous inputs form the column space of
\[K=\bigl(B\ \ AB\ \ A^2B\ \cdots\ A^{n-1}B\bigr).
\]
Consequently controllability is equivalent to $\operatorname{rank}K=n$.
\end{theorem}
\begin{proof}
The reachable set $R_T$ is a linear subspace. Cayley--Hamilton makes every $A^jB$ have columns in $\operatorname{col}K$. This finite-dimensional subspace is closed, so the exponential series and then integration show $R_T\subseteq\operatorname{col}K$.

Take $v\perp R_T$ and choose the continuous input $u(s)=B^*e^{(T-s)A^*}v$. Its endpoint gives
\[0=v^*x(T)=\int_0^T\|B^*e^{(T-s)A^*}v\|^2\dd s.
\]
The integrand is continuous and nonnegative, hence identically zero. Thus $v^*e^{\tau A}B=0$ for $0\le\tau\le T$. Derivatives at zero give $v^*A^jB=0$, so $v\perp\operatorname{col}K$. Equality of the orthogonal complements proves equality of the subspaces. For real matrices the same proof works over $\R$.
\end{proof}

The proof suggests an explicit input. Define the \emph{controllability Gramian}
\[W_T=\int_0^T e^{\tau A}BB^*e^{\tau A^*}\dd\tau.
\]
Full rank makes it positive definite. Then $u(s)=B^*e^{(T-s)A^*}W_T^{-1}x_T$ reaches $x_T$ with least energy $x_T^*W_T^{-1}x_T$; Exercise~\ref{linear:ex:mininput} proves both assertions. Reachability is independent of the positive time allowed. The least energy is non-increasing with that time: $W_T$ increases in Loewner order, and its inverse decreases.

\begin{example}[the double integrator]\label{linear:exa:doubleintegrator}
A unit-mass cart with position $q$ and force $u$ obeys $q''=u$. For the state $(q,q')$,
\[A=\begin{pmatrix}0&1\\0&0\end{pmatrix},\qquad B=\binom01,\qquad K=\begin{pmatrix}0&1\\1&0\end{pmatrix}.
\]
The rank is two: one force controls position and velocity. Two identical carts driven by the same force instead have rank two in a four-dimensional state space. Their position difference has zero second derivative regardless of the input, so an arbitrary relative displacement cannot be corrected.
\end{example}

\section{A gem: Maxwell's cubic}

Maxwell's governor model led to a real cubic. Testing its roots needs just four inequalities.
\begin{theorem}[Maxwell, 1868]\label{linear:thm:maxwell}
The roots of $s^3+as^2+bs+c$, with real coefficients, all have negative real parts exactly when $a,b,c>0$ and $ab>c$.
\end{theorem}
\begin{proof}
Factor at a real root as $(s-r)(s^2+ps+q)$. Coefficient comparison gives $a=p-r$, $b=q-rp$, $c=-rq$, and
\[ab-c=p(b+r^2).
\]
A real quadratic $s^2+ps+q$ has both roots strictly in the left half-plane exactly when $p,q>0$: real roots have negative sum and positive product, while conjugate roots have real part $-p/2$ and product equal to their squared modulus.

If the cubic is stable, $r<0$ and $p,q>0$, making $a,b,c$ and $ab-c$ positive. Conversely assume the four inequalities. The identity gives $p>0$. Since $c=-rq>0$, neither $r$ nor $q$ is zero and their signs are opposite. A positive $r$ would give $q<0$ and $b=q-rp<0$. Therefore $r<0$ and $q>0$, proving stability of both factors.
\end{proof}

The general Routh--Hurwitz criterion for a monic real polynomial requires positivity of the leading principal minors of its Hurwitz matrix. For this cubic they are $a$, $ab-c$ and $c(ab-c)$, which give Maxwell's test. We quote the higher-degree criterion; its proof needs more than this factorisation.

\begin{beauty}
An eigenvalue predicts eventual growth or decay, but a tilted pair of eigenvectors can hide a large early excursion. A quadratic energy makes that excursion easier to understand. For an algorithm we can choose an error polynomial; for a cart we can choose a force. In both cases the matrix tells us which changes are possible and what they cost.
\end{beauty}

\begin{figure}[ht]
\centering
\fig{appF-rw-energy}
\caption{For $A=\left(\begin{smallmatrix}-1&2\\0&-1\end{smallmatrix}\right)$, the energy with $P=\left(\begin{smallmatrix}1/2&1/2\\1/2&3/2\end{smallmatrix}\right)$ decreases at rate $-\|x\|_2^2$. The ellipses are its level sets; the curve starts at $(0,3/5)$.}
\label{fig:v8appFenergy}
\end{figure}
\begin{figure}[ht]
\centering
\fig{appF-rw-cart}
\caption{The least-energy trip from rest at position zero to rest at position one: $q=3(t/T)^2-2(t/T)^3$ and $T^2u=6-12t/T$. The energy cost is $12/T^3$.}
\label{fig:v8appFcart}
\end{figure}

\section{Exercises}

\begin{exercise}[$\star$]\label{linear:ex:gelfandex}
Take $A=\left(\begin{smallmatrix}1&1\\0&1\end{smallmatrix}\right)$. Find its maximum-row-sum norm $\|A^k\|_\infty$ and verify Gelfand's formula. Give matrices of spectral radius one with bounded and unbounded powers. Prove that diagonalisation and $\rho(A)\le1$ together guarantee bounded powers.
\end{exercise}
\begin{exercise}[$\star$]\label{linear:ex:hurwitz2}
Prove the planar stability test: for a real $2\times2$ matrix, Hurwitz stability is equivalent to $\operatorname{tr}A<0$ and $\det A>0$.
\end{exercise}
\begin{exercise}[$\star\star$]\label{linear:ex:stein}
Establish the discrete Lyapunov criterion: $\rho(A)<1$ exactly when, for some positive definite $Q$, there is a positive definite $P$ satisfying $P-A^*PA=Q$. Prove that then every $Q\succ0$ works and that its solution is $P=\sum_{k\ge0}(A^*)^kQA^k$.
\end{exercise}
\begin{exercise}[$\star$]\label{linear:ex:oscillator}
Write $q''+2\zeta q'+q=0$, with $\zeta>0$, as a first-order system on $(q,q')$. Prove that $E=(q^2+q'^2)/2$ never increases although $A^{\mathsf T}+A$ is only negative semidefinite. Solve $A^{\mathsf T}P+PA=-I$ and verify $P\succ0$.
\end{exercise}
\begin{exercise}[$\star$]\label{linear:ex:gdstep}
For $0<m\le L$, minimise $\max_{[m,L]}|1-\eta\lambda|$ over real $\eta$. Show that the optimal step is $2/(m+L)$ and its factor is $(\kappa-1)/(\kappa+1)$, while the step $1/L$ gives factor $1-1/\kappa$.
\end{exercise}
\begin{exercise}[$\star\star$]\label{linear:ex:chebyshevrec}
Start from $T_0=1$, $T_1=x$ and $T_{k+1}=2xT_k-T_{k-1}$. Derive $T_k(\cos\theta)=\cos(k\theta)$ and $T_k(\cosh u)=\cosh(ku)$. Determine the leading coefficient for each $k\ge1$.
\end{exercise}
\begin{exercise}[$\star\star$]\label{linear:ex:heavyjordan}
For $\kappa>1$, show that the heavy-ball endpoint matrix $M_m$ has repeated eigenvalue $\sqrt\beta$ and is not diagonalisable. Prove that its powers have norm bounded above and below by positive constant multiples of $k\beta^{k/2}$ for large $k$. Treat $\kappa=1$ separately: the matrix is still not diagonalisable, but its square is zero.
\end{exercise}
\begin{exercise}[$\star$]\label{linear:ex:kalmanosc}
Apply Kalman's test to $q''+q=u$, then to $x'=-x+u$, $y'=-2y+u$, and finally to $x'=-x+u$, $y'=-y+u$. Which systems can reach every state?
\end{exercise}
\begin{exercise}[$\star\star$]\label{linear:ex:mininput}
Assume $\operatorname{rank}K=n$. Prove $W_T\succ0$ and verify that $u_\star(s)=B^*e^{(T-s)A^*}W_T^{-1}x_T$ reaches $x_T$ from zero. Show that every other continuous input with this endpoint satisfies
\[\int_0^T\|u\|^2\dd s\ge\int_0^T\|u_\star\|^2\dd s=x_T^*W_T^{-1}x_T.
\]
\end{exercise}
\begin{exercise}[$\star$]\label{linear:ex:routh}
Decide stability of $s^3+s^2+2s+8$ and $s^3+3s^2+3s+1$. Determine all real $k$ making $s^3+2s^2+s+k$ stable, and factor the polynomial at the upper boundary.
\end{exercise}
\begin{exercise}[$\star\star$]\label{linear:ex:lognorm}
Set $\mu(A)=\lambda_{\max}((A+A^*)/2)$. Prove $\|e^{tA}\|_2\le e^{t\mu(A)}$ for $t\ge0$ and deduce stability when $\mu(A)<0$. Find a stable matrix with $\mu(A)>0$ and exhibit an initial state whose norm initially increases.
\end{exercise}
\begin{exercise}[$\star\star\star$]\label{linear:ex:leontief}
Suppose $A$ has nonnegative entries. Prove $\rho(A)<1$ exactly when $I-A$ is invertible and its inverse is entrywise nonnegative; then identify that inverse as $\sum_{k\ge0}A^k$. This condition expresses productivity in Leontief's input--output model.
\end{exercise}

\begin{exercise}[$\star\star$]\label{ex:v8cFa}
For $A=\left(\begin{smallmatrix}1/2&1\\0&1/2\end{smallmatrix}\right)$, compute $A^k$ and prove
\[\|A^k\|_2=2^{-k}\bigl(\sqrt{1+k^2}+k\bigr).
\]
Determine the integer $k\ge0$ giving the largest norm, using exact inequalities.
\end{exercise}
\begin{exercise}[$\star\star$]\label{ex:v8cFb}
Let $A_c=\left(\begin{smallmatrix}-1&c\\0&-1\end{smallmatrix}\right)$ with real $c$. Solve $A_c^{\mathsf T}P+PA_c=-I$ explicitly and prove $P\succ0$. For $|c|>2$, explain how the quadratic energy decreases strictly along every nonzero trajectory while some Euclidean norms initially increase.
\end{exercise}
\begin{exercise}[$\star\star$]\label{ex:v8cFc}
Suppose $A^*P+PA=-Q$ with $P,Q\succ0$. Prove that $A+E$ remains stable whenever
\[\|E\|_2<\frac{\lambda_{\min}(Q)}{2\|P\|_2}.
\]
Thus the same quadratic form certifies stability throughout an explicit neighbourhood of $A$.
\end{exercise}
\begin{exercise}[$\star\star$]\label{ex:v8cFd}
A cart satisfies $q''=u$ and starts at rest at zero. Among continuous forces reaching $q(T)=1$, $q'(T)=0$, find the least-energy force and trajectory. Prove that the minimum energy is $12/T^3$ and explain the curves in Figure~\ref{fig:v8appFcart}.
\end{exercise}

\section{Hints}
\begin{hint}{linear:ex:gelfandex}
The nilpotent off-diagonal part has square zero. Compare with the identity matrix for the bounded example.
\end{hint}
\begin{hint}{linear:ex:hurwitz2}
Apply the quadratic root test from Maxwell's proof to the characteristic polynomial.
\end{hint}
\begin{hint}{linear:ex:stein}
Sum the discrete energies and compare successive values of $x^*Px$. Iteration also proves uniqueness.
\end{hint}
\begin{hint}{linear:ex:oscillator}
Differentiate $E$ along the equation. For $P=\left(\begin{smallmatrix}p&r\\r&s\end{smallmatrix}\right)$, coefficient comparison gives three scalar equations.
\end{hint}
\begin{hint}{linear:ex:gdstep}
Convexity in $\lambda$ leaves only the endpoints; their optimal values have opposite signs and equal modulus.
\end{hint}
\begin{hint}{linear:ex:chebyshevrec}
Use the addition formulas for cosine and hyperbolic cosine, then compare highest powers of $x$.
\end{hint}
\begin{hint}{linear:ex:heavyjordan}
The entry below the diagonal excludes a scalar matrix. Expand powers of $qI+N$ with $N^2=0$.
\end{hint}
\begin{hint}{linear:ex:kalmanosc}
Each state space is two-dimensional, so only $B$ and $AB$ enter the test.
\end{hint}
\begin{hint}{linear:ex:mininput}
The Gramian is positive by the orthogonality argument in Kalman's proof. The difference of two admissible controls is orthogonal in $L^2$ to $u_\star$.
\end{hint}
\begin{hint}{linear:ex:routh}
Substitute the coefficients into Maxwell's four inequalities, retaining strictness at both boundaries.
\end{hint}
\begin{hint}{linear:ex:lognorm}
Bound the derivative of $\|x(t)\|^2$. A stable upper triangular matrix with a large off-diagonal entry gives the requested example.
\end{hint}
\begin{hint}{linear:ex:leontief}
For the reverse implication, apply Volume~3's Perron--Frobenius theorem to $A^{\mathsf T}$ and multiply its nonnegative eigenvector by the inverse.
\end{hint}
\begin{hint}{ex:v8cFa}
Find the larger eigenvalue of $(A^k)^*A^k$. For $k\ge3$, use $\sqrt{1+k^2}\le k+1$.
\end{hint}
\begin{hint}{ex:v8cFb}
Solve the three entries in the Lyapunov equation. Compare its energy derivative with the largest eigenvalue of the symmetric part of $A_c$.
\end{hint}
\begin{hint}{ex:v8cFc}
Bound $x^*(E^*P+PE)x$ in absolute value by $2\|P\|_2\|E\|_2\|x\|_2^2$.
\end{hint}
\begin{hint}{ex:v8cFd}
Use the double-integrator Gramian and the least-energy formula. Integrate the resulting affine force twice.
\end{hint}

\section{Solutions}

\begin{solution}{linear:ex:gelfandex}
\idea A polynomial factor in $k$ disappears under the $k$th root.
\writeup Squaring the off-diagonal matrix gives zero, hence $A^k=\left(\begin{smallmatrix}1&k\\0&1\end{smallmatrix}\right)$. Its row-sum norm is $k+1$, whose $k$th root tends to $1=\rho(A)$. These powers are unbounded, whereas the identity has the same radius and bounded powers. If $A=SDS^{-1}$ and $D$ is diagonal with entries of modulus at most one, then $\|A^k\|_2\le\|S\|_2\|S^{-1}\|_2$ for every $k$; norm equivalence gives boundedness in any fixed norm.
\end{solution}
\begin{solution}{linear:ex:hurwitz2}
\idea Trace and determinant are the two coefficients of the characteristic polynomial.
\writeup That polynomial is $s^2-(\operatorname{tr}A)s+\det A$. The real quadratic test proved above requires its linear and constant coefficients to be positive. These are exactly $-\operatorname{tr}A$ and $\det A$.
\end{solution}
\begin{solution}{linear:ex:stein}
\idea The energy lost in one step supplies the series for the total energy.
\writeup When $\rho(A)<1$, choose $r$ with $\rho(A)<r<1$. The power bound gives $\|(A^*)^kQA^k\|_2\le C^2\|Q\|_2r^{2k}$. The series converges to $P\succeq Q\succ0$, and shifting its index gives $A^*PA=P-Q$.

Conversely suppose the equation holds with $P,Q\succ0$. Since $Q=P-A^*PA\preceq P$, the number $\theta=\lambda_{\min}(Q)/\lambda_{\max}(P)$ lies in $(0,1]$. For $x_{k+1}=Ax_k$, its energy satisfies
\[x_{k+1}^*Px_{k+1}=x_k^*Px_k-x_k^*Qx_k\le(1-\theta)x_k^*Px_k.
\]
This forces every trajectory to zero, hence $\rho(A)<1$. Finally, two solutions differ by $R=A^*RA$; iterating gives $R=(A^*)^kRA^k\to0$. Thus the series is the unique solution, for each positive definite $Q$.
\end{solution}
\begin{solution}{linear:ex:oscillator}
\idea The physical energy misses instantaneous decay at zero velocity; a cross term repairs it.
\writeup The matrix is $A=\left(\begin{smallmatrix}0&1\\-1&-2\zeta\end{smallmatrix}\right)$. Along its trajectories $E'=-2\zeta q'^2\le0$, while $A^{\mathsf T}+A=\operatorname{diag}(0,-4\zeta)$ is not negative definite. Substitution into the Lyapunov equation gives
\[P=\begin{pmatrix}\zeta+\frac1{2\zeta}&\frac12\\[2pt]\frac12&\frac1{2\zeta}\end{pmatrix}.
\]
Its first diagonal entry is positive and its determinant is $1/4+1/(4\zeta^2)>0$, proving positive definiteness.
\end{solution}
\begin{solution}{linear:ex:gdstep}
\idea Equalise the two endpoint errors.
\writeup Convexity gives the maximum as $\max(|1-\eta m|,|1-\eta L|)$. For $m<L$, on $[1/L,1/m]$ the first term decreases and the second increases, so their maximum is minimised where $1-\eta m=\eta L-1$. Outside this interval moving towards it decreases the maximum. The optimum is therefore $\eta=2/(m+L)$ with value $(L-m)/(L+m)$. At $m=L$ the unique optimal step is $1/m$, giving zero as the same formulas say. The choice $1/L$ has endpoint maximum $1-m/L=1-1/\kappa$.
\end{solution}
\begin{solution}{linear:ex:chebyshevrec}
\idea The recurrence copies the addition identities.
\writeup The first two cosine identities hold directly. If they hold through index $k$, then
\[T_{k+1}(\cos\theta)=2\cos\theta\cos(k\theta)-\cos((k-1)\theta)=\cos((k+1)\theta).
\]
The hyperbolic addition identity gives the identical induction for $\cosh$. The leading term of $2xT_k$ cannot cancel against $T_{k-1}$; starting at $T_1=x$, its coefficient doubles at each subsequent step and is $2^{k-1}$.
\end{solution}
\begin{solution}{linear:ex:heavyjordan}
\idea A nontrivial Jordan block adds a factor of $k$ only when its eigenvalue is nonzero.
\writeup Put $q=\sqrt\beta>0$. At the lower spectral endpoint the characteristic polynomial is $(z-q)^2$. Since the $(2,1)$ entry is one, $M_m\ne qI$, so it is similar to $qI+N$ with $N\ne0$, $N^2=0$. Its powers are similar to $q^kI+kq^{k-1}N$. After division by $kq^k$ these matrices converge to a nonzero matrix, proving the asserted two-sided norm bounds. For $\kappa=1$, $M_m=\left(\begin{smallmatrix}0&0\\1&0\end{smallmatrix}\right)$ is nonzero with square zero. A diagonalisable matrix with sole eigenvalue zero would itself be zero, so this case too is not diagonalisable, but has no large-$k$ polynomial-times-positive-rate asymptotic.
\end{solution}
\begin{solution}{linear:ex:kalmanosc}
\idea Different responses to one input can still separate two modes.
\writeup The oscillator has $A=\left(\begin{smallmatrix}0&1\\-1&0\end{smallmatrix}\right)$ and $B=(0,1)^{\mathsf T}$, giving $K=\left(\begin{smallmatrix}0&1\\1&0\end{smallmatrix}\right)$ of rank two. For the distinct decay rates, $A=\operatorname{diag}(-1,-2)$ and $B=(1,1)^{\mathsf T}$, so $K$ has columns $(1,1)^{\mathsf T}$ and $(-1,-2)^{\mathsf T}$ and determinant $-1$. Both systems are controllable. For equal decay rates $AB=-B$, giving rank one. In that case $(x-y)'=-(x-y)$ regardless of the control, so zero initial difference stays zero.
\end{solution}
\begin{solution}{linear:ex:mininput}
\idea A minimal control is orthogonal to every input that leaves the endpoint unchanged.
\writeup For any $v$,
\[v^*W_Tv=\int_0^T\|B^*e^{\tau A^*}v\|^2\dd\tau.
\]
If this vanishes, continuity and differentiation at zero give $v\perp\operatorname{col}K$, hence $v=0$. Thus $W_T\succ0$. Set $w=W_T^{-1}x_T$. The proposed control reaches $W_Tw=x_T$. Another admissible input differs by $d$ with $\int_0^T e^{(T-s)A}Bd(s)\dd s=0$, whence
\[\int_0^T u_\star(s)^*d(s)\dd s=w^*\int_0^T e^{(T-s)A}Bd(s)\dd s=0.
\]
Expanding the squared norm gives $\int\|u\|^2=\int\|u_\star\|^2+\int\|d\|^2$. The first term is $w^*W_Tw=x_T^*W_T^{-1}x_T$, proving the minimum and its value.
\end{solution}
\begin{solution}{linear:ex:routh}
\idea The strict product inequality detects the instability.
\writeup For the first cubic $ab=2<8=c$, so it is not stable. In fact its unique real root lies below $-1$: its derivative $3s^2+2s+2$ is positive everywhere, and its value at $-1$ is six. The other two roots are conjugate and their sum $-1-r$ is positive, so both have positive real part. The second cubic is $(s+1)^3$ and is stable. The parameter family requires $k>0$ and $k<2$, hence exactly $0<k<2$. At $k=2$ it factors as $(s+2)(s^2+1)$ and has an undamped oscillatory pair; at $k=0$ it has a zero root.
\end{solution}
\begin{solution}{linear:ex:lognorm}
\idea The symmetric part measures instantaneous norm change.
\writeup Along $x'=Ax$, the derivative of $\|x\|_2^2$ is $x^*(A+A^*)x\le2\mu(A)\|x\|_2^2$. Multiplying by $e^{-2\mu(A)t}$ and integrating yields the operator bound. Negative $\mu(A)$ forces every trajectory to decay. For $A=\left(\begin{smallmatrix}-1&10\\0&-1\end{smallmatrix}\right)$, both eigenvalues are $-1$ but the symmetric part has eigenvalues $4,-6$. Starting from $x_0=(1,1)^{\mathsf T}/\sqrt2$ gives $\left.\frac{\dd}{\dd t}\|e^{tA}x_0\|_2\right|_{t=0}=4$. This is an initial derivative, not an exact $e^{4t}$ trajectory. For normal $A$, unitary diagonalisation gives $\mu(A)=\alpha(A)$; a stable normal matrix has no such increase.
\end{solution}
\begin{solution}{linear:ex:leontief}
\idea A Perron eigenvector tests the sign of $1-\rho(A)$.
\writeup If $\rho(A)<1$, the power estimate makes $M=\sum_{k\ge0}A^k$ converge. Telescoping gives $(I-A)M=M(I-A)=I$, and each entry of $M$ is nonnegative. Conversely suppose $M=(I-A)^{-1}\ge0$. Perron--Frobenius in Volume~3 gives $y\ge0$, $y\ne0$ with $A^{\mathsf T}y=\rho y$, where $\rho=\rho(A)$. Invertibility excludes $\rho=1$, and the eigenvector equation gives $y=(1-\rho)M^{\mathsf T}y$. The vector $M^{\mathsf T}y$ is nonnegative and nonzero, since $M^{\mathsf T}$ is invertible. A negative coefficient would contradict $y\ge0$, $y\ne0$. Thus $1-\rho>0$.
\end{solution}

\begin{solution}{ex:v8cFa}
\idea Compute the singular value exactly, then bound the tail of the integer sequence.
\writeup Nilpotence gives $A^k=2^{-k}\left(\begin{smallmatrix}1&2k\\0&1\end{smallmatrix}\right)$, including $k=0$. The larger eigenvalue of the latter matrix's Gram matrix is $1+2k^2+2k\sqrt{1+k^2}=(\sqrt{1+k^2}+k)^2$, proving the norm formula. At $k=1$ the norm is $(\sqrt2+1)/2>1$, exceeding $k=0$. At $k=2$, $\sqrt5+2<2\sqrt2+2$, so its norm is smaller than the value at one. For $k\ge3$, the norm is at most $(2k+1)/2^k$. This sequence decreases, since $2k+3<2(2k+1)$ for $k\ge1$, and its value at three is $7/8<1$. Therefore the unique largest norm occurs at $k=1$.
\end{solution}
\begin{solution}{ex:v8cFb}
\idea A tilted energy ellipse can contract while a round sphere expands locally.
\writeup Entry comparison gives
\[P=\begin{pmatrix}1/2&c/4\\c/4&1/2+c^2/4\end{pmatrix},\qquad \det P=\frac14+\frac{c^2}{16}>0.
\]
Its positive first entry proves $P\succ0$. The equation gives $(x^{\mathsf T}Px)'=-\|x\|_2^2$, strictly negative for a nonzero trajectory. The symmetric part of $A_c$ has largest eigenvalue $-1+|c|/2>0$ when $|c|>2$. A unit eigenvector for it has positive initial norm derivative, despite the decreasing quadratic energy.
\end{solution}
\begin{solution}{ex:v8cFc}
\idea Keep the strict negative quadratic bound after adding the perturbation.
\writeup For every $x$, Cauchy--Schwarz and the operator norms give
\[x^*\bigl((A+E)^*P+P(A+E)\bigr)x\le-\bigl(\lambda_{\min}(Q)-2\|P\|_2\|E\|_2\bigr)\|x\|_2^2.
\]
The parenthesis is positive under the stated condition. Thus the perturbed Lyapunov inequality is strictly negative definite, and Theorem~\ref{linear:thm:lyapunov} proves stability.
\end{solution}
\begin{solution}{ex:v8cFd}
\idea The Gramian determines the affine force and its exact cost.
\writeup Since $e^{\tau A}B=(\tau,1)^{\mathsf T}$, integration gives
\[W_T=\begin{pmatrix}T^3/3&T^2/2\\T^2/2&T\end{pmatrix},\qquad W_T^{-1}=\begin{pmatrix}12/T^3&-6/T^2\\-6/T^2&4/T\end{pmatrix}.
\]
For target $(1,0)^{\mathsf T}$, the least-energy formula yields $u_\star(s)=6/T^2-12s/T^3$. Twice integrating from zero gives $q(s)=3(s/T)^2-2(s/T)^3$, which has the required terminal position and velocity. The minimum cost is the first diagonal entry of $W_T^{-1}$, namely $12/T^3$. These are the normalised curves drawn in the figure.
\end{solution}

\section*{Sources and further reading}

\begin{itemize}
\item J.~C.~Maxwell, \emph{On governors}, Proceedings of the Royal Society of London 16 (1868), 270--283; E.~J.~Routh, \emph{A Treatise on
the Stability of a Given State of Motion}, Macmillan, 1877; A.~Hurwitz, \emph{\"Uber die Bedingungen, unter welchen eine Gleichung nur
Wurzeln mit negativen reellen Theilen besitzt}, Mathematische Annalen 46 (1895), 273--284.
\item A.~M.~Lyapunov, \emph{The General Problem of the Stability of Motion} (1892), English translation, Taylor \& Francis, 1992.
\item I.~Gelfand, \emph{Normierte Ringe}, Matematicheskii Sbornik 9 (1941), 3--24.
\item L.~N.~Trefethen and M.~Embree, \emph{Spectra and Pseudospectra: The Behavior of Nonnormal Matrices and Operators}, Princeton
University Press, 2005; L.~N.~Trefethen, A.~E.~Trefethen, S.~C.~Reddy and T.~A.~Driscoll, \emph{Hydrodynamic stability without
eigenvalues}, Science 261 (1993), 578--584.
\item B.~T.~Polyak, \emph{Some methods of speeding up the convergence of iteration methods}, USSR Computational Mathematics and
Mathematical Physics 4 (1964), 1--17; Y.~Nesterov, \emph{A method of solving a convex programming problem with convergence rate
$O(1/k^2)$}, Soviet Mathematics Doklady 27 (1983), 372--376.
\item M.~R.~Hestenes and E.~Stiefel, \emph{Methods of conjugate gradients for solving linear systems}, Journal of Research of the National
Bureau of Standards 49 (1952), 409--436; Y.~Saad, \emph{Iterative Methods for Sparse Linear Systems}, 2nd ed., SIAM, 2003.
\item R.~E.~Kalman, \emph{On the general theory of control systems}, Proceedings of the First IFAC Congress, Moscow, 1960; R.~E.~Kalman,
\emph{A new approach to linear filtering and prediction problems}, Journal of Basic Engineering 82 (1960), 35--45; E.~D.~Sontag,
\emph{Mathematical Control Theory}, 2nd ed., Springer, 1998.
\item O.~Mayr, \emph{The Origins of Feedback Control}, MIT Press, 1970 (Watt's governor and its ancestors).
\end{itemize}
